\section{Packing signed intervals}
\label{sec:packing}

The stationary contributions of the quadratic phases used below must
lie in disjoint intervals of the angle variable.  The interval centers
cannot be chosen separately: each group
of centers comes from evaluating several integer linear forms at one
point of a torus.  The following lemma supplies the required packing
under pairwise nonparallelism, even when larger collections of the forms
satisfy linear relations.  We give the argument from
\cite[Section ``Packing signed intervals'']{OpenAIMinimal}.

For $c\in\T$ and $0<\ell<1$, let
\[
 I(c,\ell)=\{c+x\pmod 1:-\ell/2\le x\le\ell/2\}.
\]
Thus $I(c,\ell)$ is a closed interval of length $\ell$.

\begin{lemma}[Signed interval packing]
\label{lem:signed-packing}
Let $m\ge1$ and $r\ge2$ be integers, and let
$a_1,\ldots,a_r\in\Z^m\setminus\{0\}$ be pairwise nonparallel.
If $w_1,\ldots,w_r>0$ and $2\sum_{i=1}^r w_i<1$, there are an integer
$H\ge1$ and points $\theta_1,\ldots,\theta_H\in\T^m$ for which the
$2rH$ closed intervals
\[
 I\bigl(\sigma a_i\cdot\theta_h,w_i/H\bigr),
 \qquad 1\le i\le r,\quad 1\le h\le H,\quad \sigma\in\{-1,1\},
\]
are pairwise disjoint.
\end{lemma}

We encode compatible choices of centers by a hypergraph.  Recall that
an $r$-uniform hypergraph consists of a finite vertex set and a family
of distinct $r$-element subsets, called edges.  Its degree $d(v)$ counts
edges containing $v$, and its codegree $d(u,v)$ counts edges containing
both distinct vertices $u,v$.  A matching is a family of pairwise
disjoint edges.  The combinatorial input is the following form of the
Pippenger--Spencer theorem \cite{PippengerSpencer1989}, stated with these
quantifiers in \cite[Lemma~2.1]{AlonYuster2005}.

\begin{theorem}[Pippenger--Spencer]
\label{thm:packing-edge-coloring}
For each integer $r\ge2$ and each $\gamma>0$, there is $\beta>0$ such
that the following holds.  If an $r$-uniform hypergraph satisfies, for
some $D>0$,
\[
 (1-\beta)D<d(v)<(1+\beta)D\quad\text{for every vertex }v,
 \qquad d(u,v)<\beta D\quad\text{for }u\ne v,
\]
then its edges can be partitioned into at most $(1+\gamma)D$ matchings.
\end{theorem}

In particular, a sequence of $r$-uniform hypergraphs with all degrees
$(1+o(1))D$ and maximum codegree $o(D)$ has a matching covering a
$1-o(1)$ proportion of its vertices.  Indeed, degree summation gives
$(1+o(1))|V|D/r$ edges, and the theorem partitions them into
$(1+o(1))D$ matchings; a largest matching therefore has at least
$(1-o(1))|V|/r$ edges.

\begin{proof}[Proof of Lemma~\ref{lem:signed-packing}]
We place small disjoint slots on a finite grid in a half-circle.  Each
slot has one of the types $1,\ldots,r$ and reserves both itself and its
reflection.  A compatible choice of one slot of every type will be an
edge.  We randomize the slots so that almost all vertices have the same
degree, then use Theorem~\ref{thm:packing-edge-coloring} to select many
compatible choices without reusing a slot.

\paragraph{The grid and the random slots.}
Choose positive numbers $\alpha_i>2w_i$ with $\sum_i\alpha_i=1$.
Approximating large multiples of the $\alpha_i$ by positive odd integers
gives fixed odd integers $L_i$ such that, with $L=\sum_i L_i$,
\begin{equation}
 w_i<\frac{L_i}{2L}\qquad(1\le i\le r).
 \label{eq:packing-width-slack}
\end{equation}
On a grid with spacing $1/q$, a type-$i$ slot will have length $L_i/q$.
We seek a matching of size $H\sim q/(2L)$, so that
\eqref{eq:packing-width-slack} makes the desired width $w_i/H$ strictly
smaller than the slot length.

Let $b$ be the rank over $\mathbb Q$ of the matrix with rows $a_i$.
Then $b\ge2$.  For all sufficiently large primes $q$, this matrix has
rank $b$ over $\mathbb F_q$, and every pair of its rows has rank two.
To see this, exclude the prime divisors of one nonzero $b$-by-$b$ minor
and one nonzero two-by-two minor for each pair of rows.  In the rest of
the proof $q$ tends to infinity through these primes.  All implied
constants may depend on the fixed forms and the fixed integers $L_i$.

Put $n=(q-1)/2$, $R=\lfloor\sqrt q\rfloor$, and
$U_q=\{1,\ldots,n\}$.  For $s\in\mathbb F_q$, write $|s|_q$ for the
absolute value of its representative in $[-n,n]$.  We call $|s|_q$ its
folded site.  Partition $U_q$ into consecutive chunks of $R$ sites, with
a possibly shorter final chunk; their number is $K=\lceil n/R\rceil$.

Partition the integers periodically into consecutive slots, each period
consisting of one type-$i$ slot of $L_i$ sites for each $i=1,\ldots,r$,
in that order.  In each chunk independently, translate this pattern by
a uniform residue modulo $L$ and retain only the slots contained entirely
in the chunk.  A retained slot is a vertex of its type.  Its middle site
$c$ is an integer because $L_i$ is odd.  Independently give each retained
slot the oriented center $c$ or $-c$ in $\mathbb F_q$, with equal
probabilities.  The orientation specifies which signed center will enter
the linear equations; geometrically both signs will be used.

A chunk of $u$ sites has $u/L+O(1)$ retained slots of each type, uniformly
in the shift.  Thus every realization has
\begin{equation}
 \frac{q}{2L}+O(q/R+1)
 \label{eq:packing-vertex-count}
\end{equation}
vertices of each type.  Call a site interior if its distance from both
ends of its chunk is at least $L$.  Let $C_j(s)$ indicate that $s$ is the
oriented center of a type-$j$ vertex; in particular $C_j(0)=0$.
If $|s|_q$ is interior, exactly one shift modulo $L$ puts a type-$j$
middle site there, and the required orientation has probability $1/2$.
Consequently
\begin{equation}
 \mathbb E C_j(s)=p,\qquad p=\frac1{2L},
 \label{eq:packing-center-probability}
\end{equation}
whenever $|s|_q$ is interior.  Center events in different chunks are
independent.

\paragraph{Compatibility, degrees, and codegrees.}
Declare a set of one vertex of each type, with oriented centers
$s_1,\ldots,s_r$, to be an edge when
\[
 a_i\cdot y=s_i\qquad(1\le i\le r)
\]
has a solution $y\in\mathbb F_q^m$.  Each compatible tuple is included
only once.  Since the evaluation map has rank $b$, every edge has exactly
$q^{m-b}$ generating vectors $y$.

For $s\ne0$, define
\[
 Z_{i,s}=\frac1{q^{m-1}}
 \sum_{\substack{y\in\mathbb F_q^m\\a_i\cdot y=s}}
       \prod_{j\ne i}C_j(a_j\cdot y).
\]
The affine fiber in this sum has $q^{m-1}$ elements.  If the vertex
$(i,s)$ is present, division by the multiplicity $q^{m-b}$ gives its
degree:
\begin{equation}
 d(i,s)=q^{b-1}Z_{i,s}.
 \label{eq:packing-degree}
\end{equation}
Fixing two centers of different types leaves $q^{m-2}$ possible
generating vectors, by pairwise rank two.  Their codegree is therefore
at most
\begin{equation}
 q^{b-2}.
 \label{eq:packing-codegree}
\end{equation}
Vertices of the same type have codegree zero.

We will show that almost all present vertices have degree close to
$D_q=p^{r-1}q^{b-1}$.  Joint independence of all the forms is unavailable.
Instead, most vectors in a fixed fiber put the other evaluations in
different chunks, where the slot choices are independent.  The same
observation for two fiber vectors will control the variance.

\paragraph{Conditional first and second moments.}
All relations in this paragraph are over $\mathbb F_q$.  For a fixed
type $i$, exclude the nonzero centers $s$ for which there are indices
$j<l$, both different from $i$, and a relation
\begin{equation}
 a_j+a_l=c a_i\quad\text{or}\quad a_j-a_l=c a_i,
 \qquad c\in\mathbb F_q^{\times},\qquad |cs|_q\le R.
 \label{eq:packing-exceptional-relation}
\end{equation}
Call this set $E_i$.  For either signed combination of a fixed pair,
there is at most one such $c$; if the combination is a multiple of
$a_i$, the multiple cannot be zero by pairwise nonparallelism.  Since
multiplication by $c\ne0$ is a bijection,
\begin{equation}
 |E_i|\le 2\binom{r-1}{2}(2R+1)=O(R).
 \label{eq:packing-exception-count}
\end{equation}

Fix $s\ne0$ outside $E_i$, and let $Q_s$ be the chunk containing
$|s|_q$.  Sample $y$ uniformly from $a_i\cdot y=s$, independently of
the random slots.  For each $j\ne i$, pairwise rank two makes
$a_j\cdot y$ uniform on $\mathbb F_q$.  There are at most $2LK$
noninterior positive sites, including those in a short final chunk.
Thus a specified evaluation is zero, folds to a noninterior site, or
folds into $Q_s$ with probability at most
\begin{equation}
 \tau_q=\frac{1+4LK+2R}{q}=O(q^{-1/2}).
 \label{eq:packing-boundary-probability}
\end{equation}

Two residues folding into the same chunk have either a sum or a
difference with folded absolute value at most $R$.  On the fiber
$a_i\cdot y=s$, each of $(a_j+a_l)\cdot y$ and
$(a_j-a_l)\cdot y$ is either uniform on $\mathbb F_q$ or constant.
The latter case occurs precisely when its linear form is a multiple
$c a_i$, and its value is then $cs$.  Our exclusion of $E_i$ prevents
that constant from having folded absolute value at most $R$.
The uniform case has probability at most $(2R+1)/q$.  A union bound
therefore shows that the $r-1$ evaluations fold into distinct interior
chunks, all different from $Q_s$, except on an $O(q^{-1/2})$ fraction
of the fiber.

For the second moment take independent uniform vectors $y,y'$ in the
same fiber.  The preceding argument controls collisions within each
copy.  Across copies, $a_j\cdot y$ and $a_l\cdot y'$ are independent
uniform residues whenever $j,l\ne i$, including $j=l$.  Their sum and
difference are uniform, so the same bound controls each cross-copy
collision.  For $u=1$ or $u=2$ independent fiber samples, the fraction
for which the combined list fails to occupy distinct interior chunks
outside $Q_s$ is at most
\begin{equation}
 u(r-1)\tau_q+
 2\binom{u(r-1)}2\frac{2R+1}{q}=O(q^{-1/2}).
 \label{eq:packing-good-tuples}
\end{equation}

Let $V_{i,s}$ be the event that the vertex $(i,s)$ is present, and
condition on this event whenever it has positive probability.  It
depends only on the random choices in $Q_s$.  For every fiber tuple
outside the exceptional fraction in \eqref{eq:packing-good-tuples},
the required center events occur in distinct other chunks.  Conditional
on $V_{i,s}$ they remain independent, each with probability $p$.
The exceptional fraction depends only on the fiber vectors, and its
contribution to an expectation is between zero and that fraction.
Averaging over one and two fiber vectors therefore gives
\begin{align}
 \mathbb E[Z_{i,s}\mid V_{i,s}]&=p^{r-1}+O(q^{-1/2}),
 \label{eq:packing-first-moment}\\
 \mathbb E[Z_{i,s}^2\mid V_{i,s}]&=p^{2(r-1)}+O(q^{-1/2}).
 \label{eq:packing-second-moment}
\end{align}
Here the expectations are over the slots and their orientations; the
square in the second line expands as the average over the independent
pair $y,y'$.  The bounds are uniform in $s\notin E_i$.  No interior
condition on $|s|_q$ is needed, since the entire chunk $Q_s$ was avoided.

\paragraph{Deletion and a large matching.}
Set $\varepsilon_q=q^{-1/8}$.  The two moment estimates imply
\[
 \mathbb E\bigl[(Z_{i,s}-p^{r-1})^2\mid V_{i,s}\bigr]
 =O(q^{-1/2}).
\]
By Chebyshev's inequality and \eqref{eq:packing-degree}, a present
vertex with $s\notin E_i$ has degree outside
$[(1-\varepsilon_q)D_q,(1+\varepsilon_q)D_q]$ with conditional
probability $O(q^{-1/4})$.  Summing over the at most $r(q-1)$ candidate
centers, the expected number of these vertices is $O(q^{3/4})$.
The vertices with centers in the sets $E_i$ number $O(q^{1/2})$ in
every realization.  Markov's inequality thus gives a realization with
at most $q^{7/8}$ vertices in these two classes, for sufficiently large
$q$.

Delete all these vertices and their incident edges.  A surviving vertex
loses at most the sum of its codegrees with deleted vertices, which by
\eqref{eq:packing-codegree} is at most
\[
 q^{7/8}q^{b-2}=O(q^{-1/8}D_q).
\]
Every remaining degree is consequently $(1+O(q^{-1/8}))D_q$, and the
maximum codegree is still at most $q^{b-2}=o(D_q)$.  Moreover,
\eqref{eq:packing-vertex-count} shows that each type retains
$(1+o(1))q/(2L)$ vertices.  The matching consequence of
Theorem~\ref{thm:packing-edge-coloring} gives a matching of size at least
$(1-o(1))q/(2L)$.  A matching uses one vertex of each type per edge,
so its size is also at most $(1+o(1))q/(2L)$.  We have obtained
\begin{equation}
 H=(1+o(1))\frac q{2L}.
 \label{eq:packing-matching-size}
\end{equation}

\paragraph{From slots to closed intervals.}
For each matching edge choose a generating vector
$y_h\in\mathbb F_q^m$, and set $\theta_h=y_h/q\pmod{\Z^m}$ using
integer representatives of its coordinates.  Then
$a_i\cdot\theta_h$ is the corresponding oriented center divided by
$q$, modulo one.

A slot with sites $t,\ldots,t+L_i-1$ determines the real interval
\[
 J=\left[\frac{t-1/2}{q},\frac{t+L_i-1/2}{q}\right],
\]
of length $L_i/q$ and centered at its middle site divided by $q$.
All positive slot intervals lie in $[1/(2q),1/2]$ and have disjoint
interiors.  Their reflections modulo one lie in
$[1/2,1-1/(2q)]$ and also have disjoint interiors.  The whole family
therefore has disjoint interiors, although adjacent intervals may
share endpoints, including at $1/2$.

By \eqref{eq:packing-width-slack} and
\eqref{eq:packing-matching-size}, for all sufficiently large $q$,
\[
 \frac{w_i}{H}<\frac{L_i}{q}\qquad(1\le i\le r).
\]
Each desired closed interval is thus contained strictly inside its
slot interval or its reflection.  The two signs use this same pair
regardless of the chosen orientation.  The matching never reuses a
slot, so all $2rH$ closed intervals are pairwise disjoint.
\end{proof}

Each interval in the lemma is disjoint from its reflection and hence
avoids $0$ and $1/2$.  Once the forms and widths are fixed, we fix one
prime and one matching from the proof.  The resulting finite data
$H,\theta_1,\ldots,\theta_H$ are independent of the later sampling
length $N$; in particular, this use of a finite field imposes no
arithmetic restriction on $N$.
